\chapter{热力学积分}

\section{理论框架}

热力学基本关系$S(T)=\int_{0}^{T}(C_{v}/T')\,\mathrm{d}T'+S(0)$
以微分形式表达了熵与比热之间的联系。其积分形式将$T=0$和$T=\infty$之间的熵差与
$C_{v}/T$的积分关联起来：
%
\begin{equation}
  S(\infty) - S(0) = \int_{0}^{\infty} \frac{C_{v}}{T}\,\mathrm{d}T.
  \label{eq:thermo_int}
\end{equation}
%
除以$N$得到每皇后量：
%
\begin{equation}
  \frac{S(\infty)}{N} - \frac{S(0)}{N}
  = \int_{0}^{\infty} \frac{C_{v}}{NT}\,\mathrm{d}T.
  \label{eq:DS}
\end{equation}

这一关系连接了两个独立可知的量：左边由组合分析确定，右边由Monte Carlo数据通过
数值积分确定。具体而言，左边$S(\infty)/N=(1/N)\ln\binom{N^{2}}{N}\approx\ln N+1$
是均匀分布的熵——这一公式是平凡精确的，不依赖任何关于基态N皇后解的计数知识，
只需要棋盘上有$N^{2}$个格点和$N$个皇后的组合事实。
$S(0)/N=\ln Q(N)/N\approx\ln N-\gamma$是基态熵。
右边$\int_{0}^{\infty}C_{v}/(NT)\,\mathrm{d}T$可以从Monte Carlo模拟获得的
$C_{v}(T)/N$数据通过数值积分计算。

将两边结合，可以从Monte Carlo数据反推出基态熵：
%
\begin{equation}
  s_{0}^{\rm MC} = \frac{S(\infty)}{N}
    - \int_{0}^{\infty}\frac{C_{v}}{NT}\,\mathrm{d}T,
  \label{eq:s0_MC}
\end{equation}
%
进而得到Simkin常数的热力学估计：
%
\begin{equation}
  \gamma_{\rm MC} = \ln N - s_{0}^{\rm MC}
  = \ln N - \frac{S(\infty)}{N}
    + \int_{0}^{\infty}\frac{C_{v}}{NT}\,\mathrm{d}T.
\end{equation}

关键的是，这个热力学测定方案的输入完全独立于组合学对$\gamma$的知识：唯一的输入是
均匀分布的高温熵$S(\infty)/N$——一个仅依赖棋盘几何和皇后总数的平凡量——以及
Monte Carlo产生的$C_{v}(T)/N$数据。Simkin-Nobel精确值$\gamma=1.94400(1)$
仅作为事后比较的基准，不参与积分的任何一步。这使得热力学方案构成了一条真正独立于
组合枚举的测定路径。

\section{Simkin常数的热力学测定}

数值积分的具体操作如下：我们在280个温度网格点（$T=0.05\sim400\,J$）上使用
梯形法则计算积分。为保证测得的$\gamma_{\rm MC}$误差完全可控，下面对该流程中
所有可能的误差源逐项作出解析估计。

\paragraph{(i) 低温截断误差。}在$T\ll J$区域，
比热由两能级Schottky形式主导：$C_{v}/N\sim M_{\rm eff}(J/T)^{2}\exp(-J/T)$
（见第\ref{fig:cv}图与式\ref{eq:lowT}）。因此被积函数$C_{v}/(NT)$在低温区
满足$C_{v}/(NT)\sim M_{\rm eff}\,J^{2}/T^{3}\cdot\exp(-J/T)$。
利用变量替换$u=J/T$可得截断贡献的解析上界：
%
\begin{equation}
  \int_{0}^{T_{\min}}\!\frac{C_{v}}{NT}\,\mathrm{d}T
  \;\lesssim\; M_{\rm eff}\,\Bigl(\frac{J}{T_{\min}}+1\Bigr)\exp\!\Bigl(-\frac{J}{T_{\min}}\Bigr).
\end{equation}
%
代入$T_{\min}=0.05\,J$（即$J/T_{\min}=20$）与$M_{\rm eff}\sim O(1)$：
截断贡献$\lesssim 20\cdot e^{-20}\approx 4\times 10^{-8}$。与总积分
$1+\gamma\approx 2.944$相比，相对误差量级仅$10^{-8}$，可完全忽略。

\paragraph{(ii) 高温截断误差。}在$T\gg J$区域，由高温展开
$\langle E\rangle = \langle E\rangle_{\infty} - \langle\Delta H^{2}\rangle_{\infty}/T+O(T^{-2})$
立即可得$C_{v}=\langle\Delta H^{2}\rangle_{\infty}/T^{2}+O(T^{-3})$。
对于N皇后格点气体，$\langle\Delta H^{2}\rangle_{\infty}=O(N)$
（由攻击线分解形式$H=\sum_{\alpha}\binom{n_{\alpha}}{2}$中$6N-2$条弱关联攻击线
贡献的求和给出）。因此$C_{v}/(NT)\sim O(1)/T^{3}$，截断贡献为
%
\begin{equation}
  \int_{T_{\max}}^{\infty}\!\frac{C_{v}}{NT}\,\mathrm{d}T
  \;\sim\;\frac{\langle\Delta H^{2}\rangle_{\infty}/N}{2T_{\max}^{2}}
  \;\sim\;\frac{1}{T_{\max}^{2}}.
\end{equation}
%
代入$T_{\max}=400\,J$：相对误差量级$\sim 1/(400)^{2}\approx 6\times 10^{-6}$，
亦远小于统计误差。

\paragraph{(iii) 梯形法离散化误差。}我们的温度网格在$T=0.05\sim 400\,J$上
近似按对数均匀分布，共$280$个点。对数变量$x=\ln T$下的步长
$h=\ln(T_{\max}/T_{\min})/279=\ln(8000)/279\approx 0.032$。
对数变量上的梯形法则离散化误差按经典估计为$O(h^{2})\cdot\max|f''|/12$，
其中$f(x)=T\cdot C_{v}/N\cdot e^{x}/T=C_{v}/N$（按$\mathrm{d}T=T\,\mathrm{d}x$
换元后被积函数光滑）。由此估计的相对离散化误差量级为
$h^{2}/12\approx 8.5\times 10^{-5}$，亦比$0.1\%$量级的统计误差低一个数量级以上。

\paragraph{(iv) jackknife统计误差。}如下面将看到的，对$200$个等长bin做
jackknife传播得到$\sigma(\gamma_{\rm MC})\approx 0.003$（$N=1024$，
表\ref{tab:thermo}）。这是当前误差预算中的\textbf{主导项}。

\paragraph{综合误差预算。}上述四项中，截断与离散化合计贡献小于$10^{-4}$
量级的相对误差，远在$\sigma(\gamma_{\rm MC})/\gamma\approx 1.5\times 10^{-3}$
之下，因此$\gamma_{\rm MC}$的最终不确定度由Monte Carlo采样的统计涨落
（jackknife）与N皇后渐近公式的有限$N$修正$O(N^{-\alpha})$主导，
而非数值积分本身。


统计误差的传播通过jackknife方法系统处理：对200个jackknife样本分别计算
$C_{v}(T)/N$、执行梯形积分得到200个$s_{0}^{\rm MC}$和$\gamma_{\rm MC}$的
折刀估计，这些估计值的涨落给出最终结果的标准误差。
表\ref{tab:thermo}汇总了所有$N$的热力学积分结果。

\begin{table}[htbp]
  \centering
  \caption{热力学积分得到的基态熵与Simkin常数。
  $s_{0}^{\rm MC}=S(\infty)/N-\int_{0}^{\infty}(C_{v}/NT)\,\mathrm{d}T$，
  其中$S(\infty)/N=(1/N)\ln\binom{N^{2}}{N}$是精确的。统计不确定度通过将
  jackknife误差传播通过梯形积分获得。$N\le16$与精确值$s_{0}^{\rm exact}
  =\ln Q(N)/N$比较；$N\ge32$提取$\gamma_{\rm MC}=\ln N-s_{0}^{\rm MC}$
  与精确值$\gamma=1.94400(1)$比较。}
  \label{tab:thermo}
  \begin{tabular}{ccccc}
    \toprule
    $N$ & $s_{0}^{\rm MC}$ & $s_{0}^{\rm exact}$ & $\gamma_{\rm MC}$ & 偏差 \\
    \midrule
    8    & $0.566\pm0.000$  & 0.565  & --- & 0.1\% \\
    16   & $1.031\pm0.001$  & 1.032  & --- & 0.1\% \\
    \midrule
    32   & $1.633\pm0.001$  & --- & $1.833\pm0.001$ & 5.7\% \\
    64   & $2.274\pm0.002$  & --- & $1.885\pm0.002$ & 3.0\% \\
    128  & $2.943\pm0.002$  & --- & $1.909\pm0.002$ & 1.8\% \\
    256  & $3.620\pm0.002$  & --- & $1.925\pm0.002$ & 1.0\% \\
    512  & $4.308\pm0.003$  & --- & $1.931\pm0.003$ & 0.69\% \\
    1024 & $4.985\pm0.003$  & --- & $1.946\pm0.003$ & 0.11\% \\
    \bottomrule
  \end{tabular}
\end{table}

表\ref{tab:thermo}的结果证明了热力学积分流程的可靠性以及Monte Carlo数据的精度。
在$N=8$和$N=16$处，$s_{0}^{\rm MC}$与精确基态熵的偏差仅为$0.1\%$，
验证了从模拟到积分的整个流程的准确性。对于$N\ge32$，$\gamma_{\rm MC}$从
$N=32$时的$1.833$单调收敛到$N=1024$时的$1.946\pm0.003$。
与精确组合学值$\gamma=1.94400(1)$的偏差在$N=1024$处仅为$0.11\%$。
图\ref{fig:cv}(b)展示了这一单调收敛过程。

需要指出的是，在$N=1024$处jackknife传播的统计不确定度
$\sigma(\gamma_{\rm MC})=0.003$与数值结果和精确值的偏差
$|\gamma_{\rm MC}-\gamma|\approx0.002$处于同一量级，
二者在$1\sigma$以内一致。这意味着\textbf{在单一$N=1024$处，
我们的Monte Carlo测定与组合学精确值在统计意义上无法被区分}——
这本身就构成了对$\gamma$的一个独立物理学验证：
若热力学积分流程或Monte Carlo采样存在系统偏差，本应在$\le0.1\%$精度
下显露出与$\gamma=1.94400(1)$的张力，但没有发现。
然而，仅凭单点结果尚不能严格断言残余偏差的物理来源。
下一段所讨论的尺寸标度（图\ref{fig:cv}(b)）显示$\gamma_{\rm MC}(N)$
随$N$系统性地单调收敛，提示残余差异中含有可识别的有限尺寸贡献，
即$\ln Q(N)/N=\ln N-\gamma+O(N^{-\alpha})$中次导阶的修正——
这一项才是阻碍我们以更高精度从有限$N$模拟外推到$\gamma_{\infty}$的主因。

为了尝试提取热力学极限值，我们对$N=32\sim1024$六个数据点进行了有限尺寸标度拟合：
%
\begin{equation}
  \gamma_{\rm MC}(N) = \gamma_{\infty} + a\cdot N^{-\alpha},
\end{equation}
%
拟合给出$\gamma_{\infty}=1.947\pm0.004$，$\alpha\approx0.81\pm0.06$。
这个外推值与精确值$\gamma=1.94400(1)$在$1\sigma$以内一致。需要指出，
由于$Q(N)$次渐近展开的函数形式在数学上尚未完全确定，上述三参数拟合的形式是
经验性的（\textit{ad hoc}），其外推结果应视为指示性而非确定性的。
更可靠的外推将依赖对$Q(N)$次导阶修正的进一步解析理解或扩展到更大$N$
（如$N=2048$或$4096$）的Monte Carlo数据。

综合来看，Monte Carlo模拟与热力学积分的结合为Simkin常数提供了独立于组合数学的
测定途径。该方法的核心优势有三：（1）仅依赖平凡精确的高温熵，无需$Q(N)$的任何
先验知识；（2）利用$C_{v}/N$的收敛性保证大$N$极限的存在性；
（3）Jackknife误差传播确保统计不确定度的正确估计。该方法的精度受限于N皇后公式的
有限$N$修正——随着未来将模拟扩展到更大的$N$，预计精度将进一步提升。
